\begin{recipe}
[% 
    preparationtime = {\unit[20]{min}},
    portion = {\portion{2}},
    %source = {},
    calory = {520kcal | 16g Protein | 55g KH | 24g Fett},
    price = {2,80€},
    created = {19.03.2026},
    rating = {1},
]
{Couscous-Salat mit Feta}

\graph{% pictures
    big=pic/vorspeise/couscous-salat
}

\introduction{%
Ein frischer, leichter Couscous-Salat mit knackigem Gemüse und cremigem Feta. Perfekt als Vorspeise oder als Beilage zu größeren Gerichten.
}

\ingredients{
    120 g & Couscous\index{Getreide, Pasta \& Brot!Couscous}\\
    120 ml & Gemüsebrühe\index{Vorrat \& Konserven!Gemüsebrühe}\\
    \textbf{Gemüse} & \\
    1 & Gurke\index{Gemüse \& Pilze!Gurken}\\
    2 & Tomaten\index{Gemüse \& Pilze!Tomaten}\\
    1 & Paprika\index{Gemüse \& Pilze!Paprika}\\
    \textbf{Topping} & \\
    100 g & Feta\index{Milchprodukte \& Alternativen!Feta}\\
    \textbf{Dressing} & \\
    3 EL & Olivenöl\index{Öle, Saucen \& Würzmittel!Olivenöl}\\
    2 EL & Zitronensaft\index{Obst!Zitronen}\\
    1 EL & Balsamico\index{Öle, Saucen \& Würzmittel!Balsamico}\\
    & Salz, Pfeffer\index{Kräuter \& Gewürze!Pfeffer}\index{Kräuter \& Gewürze!Salz}
}

\preparation{%
    \step%
    Couscous mit heißer Gemüsebrühe übergießen, abdecken und 5 Minuten quellen lassen. Anschließend mit einer Gabel auflockern.
    
    \step%
    Gurke, Tomaten und Paprika in kleine Würfel schneiden.
    
    \step%
    Feta zerbröseln.\\
    
    \step%
    Für das Dressing Olivenöl, Zitronensaft und Balsamico verrühren und mit Salz und Pfeffer abschmecken.
    
    \step%
    Couscous mit Gemüse und Feta vermengen, Dressing unterheben und gut durchmischen.
    
    \step%
    Kurz ziehen lassen und anschließend servieren.
}

\suggestion[Variante]{%
    Mit Oliven und getrockneten Tomaten bekommt der Salat eine deutlich intensivere, mediterrane Note.
}

\hint{
    Schmeckt am besten leicht durchgezogen nach 10–15 Minuten im Kühlschrank.
}

\end{recipe}